\documentclass[aspectratio=1610, 10pt]{beamer}
\usepackage{mobi-beamer}

% ============================================================
%  Verificación y Validación de Modelos de Simulación  (T3.2)
%  Sesiones 1 y 2 — Pregrado · MoBI · Uniandes
% ============================================================

% ── Listings con paleta MoBI ──────────────────────────────────
\usepackage{listings}
\lstset{
  basicstyle=\ttfamily\scriptsize,
  keywordstyle=\color{mobiblue}\bfseries,
  commentstyle=\color{deepgreen}\itshape,
  stringstyle=\color{red!70!black},
  numberstyle=\tiny\color{gray},
  numbers=left, stepnumber=1,
  breaklines=true,
  frame=single, framerule=0.4pt,
  rulecolor=\color{gray!50},
  backgroundcolor=\color{mobisoftgray},
  showstringspaces=false, tabsize=2
}

% ── Estilo común para pgfplots (igual a T3.1) ─────────────────
\pgfplotsset{
  every axis/.append style={
    font=\scriptsize,
    label style={font=\scriptsize},
    title style={font=\footnotesize\bfseries},
    tick label style={font=\scriptsize},
    legend style={font=\scriptsize, draw=mobiblue!30, fill=white,
                  fill opacity=0.9, text opacity=1},
    grid=both, grid style={mobiblue!10, very thin},
    major grid style={mobiblue!20},
    axis line style={mobiblue!70},
    tick style={mobiblue!70}
  }
}

\title[T3.2 V\&V]{%
  Verificación y Validación de Modelos de Simulación\\[2pt]
  \large ¿Construimos el modelo correctamente y construimos el modelo correcto?}
\subtitle{T3.2 · Modelado de Sistemas bajo Incertidumbre}
\author{Alejandra Tabares Pozos}
\institute{Universidad de los Andes · Ingeniería Industrial}
\date{2026}

% ============================================================
\begin{document}
% ============================================================

\begin{frame}\titlepage\end{frame}

% ─────────────────────────────────────────────────────────────
% AGENDA
% ─────────────────────────────────────────────────────────────
\begin{frame}{Estructura de las dos sesiones}
  \begin{columns}[T]
    \begin{column}{0.48\textwidth}
      \begin{block}{Sesión 1 — Verificación}
        \footnotesize
        \emph{¿El código implementa el modelo conceptual?}
        \begin{itemize}\itemsep1pt
          \item Distinción formal V\&V (marco de Sargent).
          \item \emph{Debugging} con aserciones e invariantes.
          \item Trazas de eventos y verificación contra cálculo manual.
          \item Pruebas unitarias (patrón AAA).
          \item Verificación contra resultado analítico ($M/M/1$).
          \item \textbf{Ejemplo 1} — verificación completa de un $M/M/1$ básico.
        \end{itemize}
      \end{block}
    \end{column}
    \begin{column}{0.48\textwidth}
      \begin{alertblock}{Sesión 2 — Validación}
        \footnotesize
        \emph{¿El modelo conceptual representa al sistema real?}
        \begin{itemize}\itemsep1pt
          \item Niveles de validación (Sargent).
          \item \emph{Face validity}: protocolo y cuestionario.
          \item Validación de supuestos: $\chi^2$ y Kolmogorov--Smirnov.
          \item Calibración y división de muestra.
          \item Validación I/O con prueba $t$ pareada.
          \item Test de Turing (Schruben).
          \item \textbf{Ejemplo 2} — bug de parametrización detectado.
          \item \textbf{Ejemplo 3} — celda $M/E_2/1$ + sensibilidad.
        \end{itemize}
      \end{alertblock}
    \end{column}
  \end{columns}

  \vspace{4pt}
  \begin{notabox}[Idea central]
    \footnotesize
    Validar y verificar no son trámites finales: son las herramientas con las que
    un modelo gana \textbf{credibilidad}. Sin V\&V, ningún tomador de decisiones
    debería actuar sobre los resultados de la simulación.
  \end{notabox}
\end{frame}

% =============================================================
\section{Distinción V\&V}
% =============================================================

\begin{frame}{Las dos preguntas: validación y verificación}
  \begin{columns}[T]
    \begin{column}{0.48\textwidth}
      \begin{alertabox}[Validación]
        \footnotesize
        \emph{¿Estamos construyendo el modelo \textbf{correcto}?}\\[2pt]
        Proceso de establecer que el \textbf{modelo conceptual} es una
        representación adecuada del sistema real \emph{para los propósitos
        del estudio}.\\[2pt]
        $\Rightarrow$ Compara \textbf{modelo} contra \textbf{realidad}.
      \end{alertabox}
    \end{column}
    \begin{column}{0.48\textwidth}
      \begin{definicionbox}[Verificación]
        \footnotesize
        \emph{¿Estamos construyendo el modelo \textbf{correctamente}?}\\[2pt]
        Proceso de establecer que la \textbf{implementación computacional}
        corresponde fielmente al \textbf{modelo conceptual}.\\[2pt]
        $\Rightarrow$ Compara \textbf{código} contra \textbf{modelo conceptual}.
      \end{definicionbox}
    \end{column}
  \end{columns}

  \vspace{6pt}
  \begin{notabox}[Orden conceptual vs.\ orden de ejecución]
    \footnotesize
    \textbf{Conceptualmente, la validación es primaria}: si el modelo conceptual no
    representa al sistema, no importa qué tan bien programado esté el código.
    \textbf{En la práctica}, se verifica el código a medida que se programa y se
    valida el modelo de manera \emph{iterativa} contra la realidad.
  \end{notabox}
\end{frame}

% ------------------------------------------------------------
\begin{frame}{Ubicación de V\&V en el ciclo de modelado (Sargent)}
  \begin{center}
    \begin{tikzpicture}[scale=0.95, transform shape, >=Stealth]
      \node[draw=mobiblue, thick, fill=mobilightblue, rounded corners,
            minimum width=3.4cm, minimum height=1cm, align=center]
            (real) at (0,0) {Sistema\\real};
      \node[draw=mobiblue, thick, fill=mobilightblue, rounded corners,
            minimum width=3.4cm, minimum height=1cm, align=center]
            (conc) at (6,3) {Modelo\\conceptual};
      \node[draw=mobiblue, thick, fill=mobilightblue, rounded corners,
            minimum width=3.4cm, minimum height=1cm, align=center]
            (comp) at (12,0) {Modelo\\computacional};

      \draw[->, very thick, mobiblue]
        (real) -- node[above left, font=\footnotesize, sloped]{Análisis y abstracción} (conc);
      \draw[->, very thick, mobiblue]
        (conc) -- node[above right, font=\footnotesize, sloped]{Programación} (comp);
      \draw[<->, very thick, princetonorange, dashed]
        (real.east) -- node[below=2pt, font=\footnotesize\bfseries]
        {Validación operacional (I/O)} (comp.west);

      \node[draw=princetonorange, thick, fill=mobilightgold,
            rounded corners, font=\footnotesize\bfseries, align=center]
            at (3,2) {VALIDACIÓN\\conceptual};
      \node[draw=mobiblue, thick, fill=mobilightblue,
            rounded corners, font=\footnotesize\bfseries, align=center]
            at (9,2) {VERIFICACIÓN};

      \node[font=\scriptsize, align=center, gray!70!black] at (3,0.4)
        {Datos del sistema};
      \node[font=\scriptsize, align=center, gray!70!black] at (9,0.4)
        {Experimentación};
    \end{tikzpicture}
  \end{center}

  \vspace{2pt}
  \footnotesize
  La \textbf{validación} actúa entre el sistema real y el modelo (conceptual y
  operacional). La \textbf{verificación} actúa exclusivamente entre el modelo
  conceptual y su implementación.\\
  \emph{Adaptado de Sargent (2013) y Banks et al.\ (2010).}
\end{frame}

% ------------------------------------------------------------
\begin{frame}{Comparación lado a lado}
  \footnotesize
  \begin{center}
    \renewcommand{\arraystretch}{1.3}
    \begin{tabular}{p{3.6cm}|p{5.1cm}|p{5.1cm}}
      \toprule
      & \textbf{\textcolor{red!70!black}{Validación}}
      & \textbf{\textcolor{mobiblue}{Verificación}} \\
      \midrule
      Pregunta            & ¿El modelo correcto?
                          & ¿Construido correctamente? \\
      \hline
      Compara             & Modelo vs.\ realidad
                          & Código vs.\ modelo conceptual \\
      \hline
      Naturaleza          & Estadística y de modelado
                          & Ingeniería de software \\
      \hline
      Errores que detecta & Supuestos inadecuados, distribuciones mal ajustadas,
                            lógica de negocio incorrecta
                          & Bugs, errores lógicos, errores numéricos \\
      \hline
      Herramientas        & Pruebas estadísticas, expertos del dominio,
                            datos históricos
                          & Debugger, trazas, pruebas unitarias \\
      \hline
      ¿Cuándo se hace?    & Antes, durante y después de modelar
                          & Durante y después de programar \\
      \bottomrule
    \end{tabular}
  \end{center}

  \vspace{6pt}
  \begin{alertabox}[Moraleja]
    \footnotesize
    Un modelo puede estar \textbf{verificado} (sin bugs) pero no
    \textbf{validado} (con supuestos incorrectos), y viceversa. Ambos procesos
    son necesarios; ninguno sustituye al otro.
  \end{alertabox}
\end{frame}

% ------------------------------------------------------------
\begin{frame}{Marco de Sargent: tres tipos de validez}
  \begin{columns}[T]
    \begin{column}{0.55\textwidth}
      \begin{tikzpicture}[>=Stealth,
        layer/.style={rectangle, rounded corners=4pt, draw=mobiblue, thick,
          fill=mobilightblue, minimum width=4.6cm, minimum height=0.9cm,
          align=center, font=\footnotesize}]
        \node[layer, fill=mobilightgreen, draw=deepgreen] (real)
          {Sistema real (operación, datos)};
        \node[layer, above=0.7cm of real] (conc)
          {Modelo conceptual (supuestos, lógica)};
        \node[layer, above=0.7cm of conc] (comp)
          {Modelo computacional (código + datos)};

        \draw[->, thick, deepgreen]
          (real.east) to[bend right=55] node[right, font=\scriptsize]
          {\textbf{V. operacional}} (comp.east);
        \draw[->, thick, mobiblue]
          (conc.west) to[bend left=35] node[left, font=\scriptsize]
          {\textbf{V. computacional}} (comp.west);
        \draw[<-, thick, princetonorange]
          (real.west) to[bend left=35] node[left, font=\scriptsize, align=right]
          {\textbf{V. conceptual}} (conc.west);
      \end{tikzpicture}
    \end{column}
    \begin{column}{0.43\textwidth}
      \footnotesize
      \begin{enumerate}\itemsep3pt
        \item \textbf{Validez conceptual}\\
          ¿Es razonable la teoría y la estructura del modelo? Supuestos,
          distribuciones (T3.1), lógica de eventos.
        \item \textbf{Validez computacional}\\
          ¿La implementación reproduce el modelo conceptual? Es la verificación.
        \item \textbf{Validez operacional}\\
          ¿Las salidas del modelo coinciden con las del sistema real
          dentro de un margen aceptable?
      \end{enumerate}
      \vspace{4pt}
      \emph{Validez de datos} cruza las tres capas.
    \end{column}
  \end{columns}

  \vspace{4pt}
  \begin{notabox}[V\&V es un proceso continuo]
    \footnotesize
    Cada modificación al modelo (nueva distribución, nuevo evento, nuevo parámetro)
    requiere re-verificación y potencialmente re-validación. Un modelo validado
    para un rango de operación puede no serlo fuera de él.
  \end{notabox}
\end{frame}

% =============================================================
\section*{}
\begin{frame}[plain,c]
  \begin{center}
    {\Huge \color{mobiblue}\textbf{Sesión 1}}\\[0.5cm]
    {\Large Verificación del modelo computacional}\\[0.3cm]
    {\large \emph{Debugging}, trazas y pruebas unitarias}\\[0.4cm]
    {\color{mobigold}\rule{0.5\textwidth}{1.2pt}}
  \end{center}
\end{frame}

% =============================================================
\section{Sesión 1: Verificación}
% =============================================================

\begin{frame}{Panorama de técnicas de verificación}
  \footnotesize
  \textbf{Cuatro categorías} (Banks et al., cap.\ 10; Law, cap.\ 5):

  \vspace{4pt}
  \begin{enumerate}\itemsep4pt
    \item \textbf{Inspección estática del código}
      \begin{itemize}\scriptsize\itemsep0pt
        \item Revisión por pares (\emph{code review}, \emph{walkthroughs}).
        \item Análisis estático automatizado (\emph{linters}).
      \end{itemize}
    \item \textbf{Inspección dinámica de la ejecución}
      \hfill {\small\textit{(\textbf{debugging} y \textbf{trazas})}}
      \begin{itemize}\scriptsize\itemsep0pt
        \item Depuradores interactivos, \emph{logging}, \emph{breakpoints}, \emph{watch}.
        \item Trazas de eventos y de estado.
      \end{itemize}
    \item \textbf{Pruebas formales}
      \hfill {\small\textit{(\textbf{pruebas unitarias})}}
      \begin{itemize}\scriptsize\itemsep0pt
        \item Pruebas unitarias (\emph{unit tests}).
        \item Pruebas de integración.
        \item Pruebas de regresión.
      \end{itemize}
    \item \textbf{Pruebas de coherencia con resultados conocidos}
      \begin{itemize}\scriptsize\itemsep0pt
        \item Comparación con soluciones analíticas ($M/M/1$, $M/M/c$, P-K).
        \item Casos límite (degeneración, continuidad).
      \end{itemize}
  \end{enumerate}

  \vspace{4pt}
  \begin{metodobox}[Filosofía operativa]
    \scriptsize
    Una aserción violada \emph{detiene} la simulación y \emph{señala} dónde está
    el error. Es radicalmente más eficiente que mirar reportes de salida tratando
    de adivinar qué celda contiene el bug.
  \end{metodobox}
\end{frame}

% =============================================================
\subsection{Debugging}
% =============================================================

\begin{frame}{\emph{Debugging}: definición y herramientas formales}
  \begin{definicionbox}[\emph{Debugging}]
    \footnotesize
    Proceso sistemático de identificar, localizar y corregir defectos
    (\emph{bugs}) en el código. Tres fases:
    \emph{(i)} \textbf{detección} (síntoma observable),
    \emph{(ii)} \textbf{localización} (causa raíz),
    \emph{(iii)} \textbf{corrección} (solución verificada).
  \end{definicionbox}

  \vspace{4pt}
  \footnotesize
  \textbf{Herramientas formales que usaremos:}
  \begin{center}\scriptsize
    \renewcommand{\arraystretch}{1.2}
    \begin{tabular}{p{3.2cm} p{4.6cm} p{6.0cm}}
      \toprule
      \textbf{Técnica} & \textbf{Mecanismo} & \textbf{Cuándo se usa} \\
      \midrule
      Aserciones (\texttt{assert})
        & Detiene la ejecución si una invariante se viola
        & Verificar pre/post-condiciones e invariantes \\
      \emph{Logging} estructurado
        & Registra eventos con marca de tiempo y nivel
        & Ejecuciones largas, reproducción a posteriori \\
      \emph{Breakpoints}
        & Pausa en líneas o condiciones específicas
        & Análisis interactivo del estado \\
      Inspección de variables (\emph{watch})
        & Vigila el cambio de un valor
        & Detectar cuándo una variable toma valor inesperado \\
      Reproducibilidad por semilla
        & Fija la semilla aleatoria
        & Aislar bugs intermitentes \\
      \bottomrule
    \end{tabular}
  \end{center}
\end{frame}

% ------------------------------------------------------------
\begin{frame}[fragile]{Debugging con aserciones — ejemplo $M/M/1$}
  \footnotesize
  \textbf{Invariantes que debe respetar el simulador:}
  \[
    N_{\text{sis}}(t) = N_q(t) + N_s(t),\quad N\ge 0,\quad
    t_{\text{ev},k+1}\ge t_{\text{ev},k}
  \]

\begin{lstlisting}[language=Python]
def manejar_salida(self, t):
    # Pre-condiciones (deben cumplirse SIEMPRE)
    assert t >= self.reloj, f"Tiempo retrocede: {t} < {self.reloj}"
    assert self.n_servicio == 1, "Salida sin cliente en servicio"
    assert self.n_cola >= 0, f"Cola negativa: {self.n_cola}"

    self.reloj = t
    self.servidos += 1
    self.n_servicio = 0

    if self.n_cola > 0:                       # Toma el siguiente
        self.n_cola -= 1
        self.n_servicio = 1
        self.programar_salida(t + self.gen_servicio())

    # Post-condicion: invariante de conservacion
    assert self.llegados == self.servidos + self.n_cola + self.n_servicio, \
           "Violacion de conservacion de clientes"
\end{lstlisting}
\end{frame}

% ------------------------------------------------------------
\begin{frame}[fragile]{\emph{Logging} estructurado: bugs sutiles}
\begin{lstlisting}[language=Python]
import logging
logging.basicConfig(level=logging.DEBUG,
                    format='%(asctime)s [%(levelname)s] %(message)s')

def manejar_llegada(self, t):
    self.reloj = t
    self.llegados += 1
    logging.debug(f"LLEGADA t={t:.4f} cliente={self.llegados} cola={self.n_cola}")

    if self.n_servicio == 0:
        self.n_servicio = 1
        ts = self.gen_servicio()
        self.programar_salida(t + ts)
        logging.info(f"INICIO SERV. t={t:.4f} servicio={ts:.4f}")
    else:
        self.n_cola += 1
        if self.n_cola > 100:
            logging.warning(f"COLA GRANDE: {self.n_cola} en t={t:.4f}")

    self.programar_llegada(t + self.gen_interarribo())
\end{lstlisting}

  \vspace{2pt}
  \begin{notabox}[Por qué funciona]
    \footnotesize
    Reconstruye la historia completa de la corrida y permite filtrar por nivel
    (\texttt{DEBUG}, \texttt{INFO}, \texttt{WARNING}). Detecta condiciones
    anómalas \emph{sin detener} la simulación.
  \end{notabox}
\end{frame}

% =============================================================
\subsection{Trazas}
% =============================================================

\begin{frame}{Trazas: definición formal}
  \begin{definicionbox}[Traza]
    \footnotesize
    Una \textbf{traza} es la secuencia ordenada de \emph{tuplas de estado} del
    sistema simulado, registradas en cada cambio de estado:
    \[
      \mathcal{T} = \big\{ (t_k, e_k, \mathbf{s}_k) : k=1,2,\ldots,K \big\}
    \]
    donde $t_k$ es el reloj de simulación, $e_k$ el evento ejecutado y
    $\mathbf{s}_k$ el vector de estado del sistema inmediatamente después
    del evento.
  \end{definicionbox}

  \vspace{6pt}
  \footnotesize\textbf{Tres usos formales de la traza para verificación:}
  \begin{enumerate}\itemsep2pt
    \item \textbf{Inspección manual} (\emph{eyeballing}): primeras 30--50
      líneas para revisar lógica.
    \item \textbf{Verificación de invariantes}: chequear automáticamente
      propiedades sobre $\mathcal{T}$ (conservación, no-negatividad, monotonía).
    \item \textbf{Comparación con cálculo a mano}: ejecutar el modelo con
      \emph{inputs} prefijados (no aleatorios) y verificar la traza
      contra resultados calculados manualmente.
  \end{enumerate}
\end{frame}

% ------------------------------------------------------------
\begin{frame}[fragile]{Traza $M/M/1$ con \emph{inputs} prefijados}
  \footnotesize
  \textbf{Inputs deterministas} (un solo cajero):
  \begin{itemize}\scriptsize\itemsep0pt
    \item Tiempos entre llegadas: $1{,}2;\;0{,}5;\;1{,}8;\;0{,}6;\;2{,}1$
    \item Tiempos de servicio: \;\;\;\;\;\;\;\;$0{,}8;\;1{,}5;\;0{,}4;\;1{,}1;\;0{,}7$
  \end{itemize}

  \vspace{4pt}
  \textbf{Traza esperada (calculada a mano):}
\begin{lstlisting}
  t        Evento    Cliente  N_cola  N_serv  Acumulado
  1.200    LLEG      C1       0       1       L=1
  2.000    SAL       C1       0       0       (servicio = 0.8)
  1.700    LLEG      C2       0       1       (1.2+0.5)
  3.200    SAL       C2       0       0       (servicio = 1.5)
  3.500    LLEG      C3       0       1       (1.7+1.8)
  3.900    SAL       C3       0       0       (servicio = 0.4)
  4.100    LLEG      C4       0       1       (3.5+0.6)
  5.200    SAL       C4       0       0       (servicio = 1.1)
  6.200    LLEG      C5       0       1
  6.900    SAL       C5       0       0
\end{lstlisting}

  \vspace{2pt}
  \begin{alertabox}[Regla de oro]
    \footnotesize
    La traza producida por el simulador debe ser \emph{idéntica} línea por línea.
    Cualquier diferencia $\Rightarrow$ \textbf{bug localizado en ese evento}.
  \end{alertabox}
\end{frame}

% ------------------------------------------------------------
\begin{frame}{Verificación automática de invariantes sobre la traza}
  \textbf{Invariantes formales} que se chequean al recorrer $\mathcal{T}$:

  \vspace{4pt}
  \footnotesize
  \begin{center}
    \renewcommand{\arraystretch}{1.25}
    \begin{tabular}{p{4.5cm} p{8.5cm}}
      \toprule
      \textbf{Invariante} & \textbf{Expresión formal} \\
      \midrule
      Monotonicidad del reloj      & $t_{k+1} \geq t_k \;\;\forall k$ \\
      No-negatividad de la cola    & $N_{q,k} \geq 0 \;\;\forall k$ \\
      Conservación de clientes     & $A_k = D_k + N_{\text{sis},k}$
                                     (arrivals = departures + WIP) \\
      Capacidad del servidor       & $0 \leq N_{\text{serv},k} \leq c$ \\
      FCFS (orden de salida)       & si $a_i < a_j \Rightarrow d_i \leq d_j$ \\
      Trabajo conservado (Little)  & $L = \lambda W$ en estado estacionario
                                     (con tolerancia) \\
      \bottomrule
    \end{tabular}
  \end{center}

  \vspace{6pt}
  \begin{metodobox}[Implementación]
    \footnotesize
    Un \emph{post-procesador} lee la traza y verifica cada invariante.
    Si alguna falla $\Rightarrow$ línea exacta donde aparece el bug.
    El costo computacional es \emph{lineal} en el número de eventos,
    despreciable frente al de la simulación.
  \end{metodobox}
\end{frame}

% =============================================================
\subsection{Pruebas unitarias}
% =============================================================

\begin{frame}{Pruebas unitarias: definición y patrón AAA}
  \begin{definicionbox}[Prueba unitaria]
    \footnotesize
    Programa que ejecuta una unidad mínima de código (función, método o módulo)
    con \emph{inputs} controlados y verifica automáticamente que el \emph{output}
    coincida con el valor esperado.
  \end{definicionbox}

  \vspace{4pt}
  \footnotesize
  \textbf{Patrón \emph{Arrange–Act–Assert} (AAA):}
  \begin{enumerate}\itemsep1pt
    \item \textbf{Arrange}: preparar el escenario (datos, semilla, instancia).
    \item \textbf{Act}: ejecutar la unidad bajo prueba.
    \item \textbf{Assert}: comparar el resultado con el valor esperado.
  \end{enumerate}

  \vspace{4pt}
  \textbf{Tres categorías de pruebas para un simulador:}
  \begin{itemize}\itemsep1pt
    \item \textbf{Pruebas de límite}: $\lambda=0$, $\mu\to\infty$, capacidad $\to 0$.
    \item \textbf{Pruebas de invariantes}: conservación, no-negatividad, FCFS.
    \item \textbf{Pruebas contra valor analítico}: $M/M/1$, $M/M/c$, $M/M/1/K$, $M/G/1$ (P-K).
  \end{itemize}
\end{frame}

% ------------------------------------------------------------
\begin{frame}[fragile]{Pruebas unitarias para un $M/M/1$ (\texttt{pytest})}
\begin{lstlisting}[language=Python]
import pytest
from simulador import SimuladorMM1

def test_sin_llegadas_servidor_ocioso():
    sim = SimuladorMM1(lam=0.0, mu=1.0, semilla=42)
    sim.correr(t_max=1000)
    assert sim.servidos == 0
    assert sim.utilizacion() == 0.0

def test_conservacion_clientes():
    sim = SimuladorMM1(lam=0.8, mu=1.0, semilla=42)
    sim.correr(t_max=1000)
    assert sim.llegados == sim.servidos + sim.en_sistema()

def test_estado_estacionario_vs_analitico():
    """L = rho/(1-rho) para M/M/1; rho=0.5 -> L=1.0"""
    sim = SimuladorMM1(lam=0.5, mu=1.0, semilla=42)
    sim.correr(t_max=200_000, t_warmup=20_000)
    assert sim.L_promedio() == pytest.approx(1.0, abs=0.05)

def test_rho_mayor_uno_cola_crece():
    sim = SimuladorMM1(lam=1.2, mu=1.0, semilla=42)
    sim.correr(t_max=10_000)
    assert sim.cola_max() > 100   # Sistema inestable
\end{lstlisting}
\end{frame}

% ------------------------------------------------------------
\begin{frame}{$M/M/1$ como banco de pruebas — verificación contra teoría}
  \footnotesize
  \textbf{Por qué $M/M/1$:} solución cerrada conocida; permite comprobar el
  simulador con alta precisión.

  \vspace{2pt}
  \begin{formulabox}[Fórmulas analíticas]
    \footnotesize
    Con $\rho=\lambda/\mu<1$:
    \[
      L=\frac{\rho}{1-\rho},\quad
      L_q=\frac{\rho^2}{1-\rho},\quad
      W=\frac{1}{\mu-\lambda},\quad
      W_q=\frac{\rho}{\mu-\lambda}
    \]
  \end{formulabox}

  \vspace{4pt}
  \textbf{Tabla de verificación} ($\lambda=0{,}8$; $\mu=1{,}0$; $\rho=0{,}8$;
  30 réplicas; $T=10\,000$; \emph{warm-up} $=1\,000$):
  \begin{center}\scriptsize
    \renewcommand{\arraystretch}{1.25}
    \begin{tabular}{lcccc}
      \toprule
      \textbf{Métrica} & \textbf{Analítico} & \textbf{Simulación}
        & \textbf{IC 95\%} & \textbf{¿Aprueba?} \\
      \midrule
      $L$    & $4{,}000$ & $4{,}018$ & $[3{,}91;\;4{,}13]$ & \textcolor{deepgreen}{$\checkmark$} \\
      $L_q$  & $3{,}200$ & $3{,}215$ & $[3{,}10;\;3{,}33]$ & \textcolor{deepgreen}{$\checkmark$} \\
      $W$    & $5{,}000$ & $5{,}027$ & $[4{,}88;\;5{,}17]$ & \textcolor{deepgreen}{$\checkmark$} \\
      $W_q$  & $4{,}000$ & $4{,}020$ & $[3{,}88;\;4{,}16]$ & \textcolor{deepgreen}{$\checkmark$} \\
      \bottomrule
    \end{tabular}
  \end{center}

  \vspace{2pt}
  \begin{notabox}[Criterio formal]
    \footnotesize
    Verificación aceptada si el valor analítico cae dentro del IC al 95\% para
    \textbf{todas} las métricas.
  \end{notabox}
\end{frame}

% ------------------------------------------------------------
\begin{frame}{$W_q$ vs $\rho$: prueba de coherencia teórica}
  \begin{columns}[T]
    \begin{column}{0.55\textwidth}
      \begin{tikzpicture}
        \begin{axis}[width=\textwidth, height=5.4cm,
          xlabel={$\rho$ (utilización)},
          ylabel={$W_q$ (unidades de tiempo)},
          xmin=0, xmax=0.99, ymin=0, ymax=20,
          legend pos=north west, samples=120]
          \addplot[thick, mobiblue, domain=0.01:0.97] {x/(1-x)};
          \addlegendentry{Teórico $M/M/1$ ($\mu=1$)}
          \addplot[only marks, mark=*, mark size=1.5pt, princetonorange]
            coordinates {
              (0.0,0.00) (0.1,0.10) (0.2,0.27) (0.3,0.42) (0.4,0.69)
              (0.5,1.05) (0.6,1.55) (0.7,2.40) (0.8,4.02) (0.9,9.18)
              (0.95,18.50)
            };
          \addlegendentry{Simulación (30 réplicas)}
        \end{axis}
      \end{tikzpicture}
    \end{column}
    \begin{column}{0.42\textwidth}
      \footnotesize
      \textbf{Lectura:}
      \begin{itemize}\itemsep1pt
        \item Los puntos siguen la curva en todo el rango.
        \item Comportamiento asintótico correcto cuando $\rho\to 1$
          ($W_q\to\infty$).
        \item No hay desviaciones sistemáticas (sesgo).
      \end{itemize}

      \vspace{4pt}
      \begin{alertabox}[Si la curva no diverge]
        \footnotesize
        Hay un truncamiento o un límite ``invisible'' en el código (cola con
        capacidad accidentalmente fijada).
      \end{alertabox}
    \end{column}
  \end{columns}
\end{frame}

% =============================================================
% EJEMPLO 1 — Verificación completa M/M/1 básico
% =============================================================
\section{Ejemplo 1 (mínimo)}

\begin{frame}{Ejemplo 1 — Verificación completa de un $M/M/1$}
  \begin{columns}[T]
    \begin{column}{0.55\textwidth}
      \begin{ejemplobox}[Sistema]
        \footnotesize
        Cola $M/M/1$ con $\lambda=3$ clientes/min y $\mu=5$ clientes/min.
        Antes de usar el modelo para análisis, se aplica el catálogo
        completo de pruebas de verificación.
      \end{ejemplobox}

      \vspace{4pt}
      \textbf{Valores analíticos de referencia:}
      \[
        \rho=\frac{\lambda}{\mu}=0{,}60
      \]
      \[
        W_q=\frac{\rho}{\mu(1-\rho)}=\frac{0{,}6}{5\times 0{,}4}=0{,}30 \text{ min}
      \]
      \[
        L_q=\frac{\rho^2}{1-\rho}=\frac{0{,}36}{0{,}4}=0{,}90 \text{ clientes}
      \]
    \end{column}
    \begin{column}{0.42\textwidth}
      \begin{tikzpicture}[>=Stealth]
        \node[circle, draw=mobiblue, fill=mobilightblue, thick,
              minimum size=0.9cm, font=\scriptsize] (q) at (0,0) {Cola};
        \node[circle, draw=deepgreen, fill=mobilightgreen, thick,
              minimum size=0.9cm, font=\scriptsize] (s) at (1.7,0) {$\mu$};
        \node[font=\scriptsize] (in) at (-1.6,0) {$\lambda=3$};
        \node[font=\scriptsize] (out) at (3.4,0) {salidas};
        \draw[->, thick, mobiblue] (in) -- (q);
        \draw[->, thick, mobiblue] (q) -- (s);
        \draw[->, thick, deepgreen] (s) -- (out);
        \node[font=\tiny] at (0.85,-0.8) {$\rho=0{,}6$};
      \end{tikzpicture}

      \vspace{4pt}
      \textbf{Pruebas a realizar:}
      \begin{itemize}\footnotesize\itemsep1pt
        \item Degenerada $\lambda=0$
        \item Conservación
        \item Traza manual
        \item Tráfico extremo $\rho\to 1$
      \end{itemize}
    \end{column}
  \end{columns}
\end{frame}

\begin{frame}{Ejemplo 1 — Resultados de las pruebas}
  \begin{center}\footnotesize
    \renewcommand{\arraystretch}{1.25}
    \begin{tabular}{lll}
      \toprule
      \textbf{Prueba} & \textbf{Esperado} & \textbf{Observado} \\
      \midrule
      Degenerada $\lambda=0$, $T=1000$
        & $N(t)\equiv 0$, $\rho=0$, $W_q=0$
        & Exactamente $0$ en todas las métricas $\checkmark$ \\
      Conservación, $T=1000$
        & $\sum_{\text{lleg}} = \sum_{\text{sal}} + N(T)$
        & $3\,012 = 3\,012 + 0$ $\checkmark$ \\
      Traza primeros 6 eventos
        & Coincide con cálculo manual
        & Tiempos exactos al 4to decimal $\checkmark$ \\
      Tráfico extremo, $\rho=0{,}998$, $T=5000$
        & $L_q$ crece sin cota
        & $L_q > 500$ al final $\checkmark$ \\
      \bottomrule
    \end{tabular}
  \end{center}

  \vspace{8pt}
  \begin{tikzpicture}
    \begin{axis}[width=12cm, height=3.4cm,
      xlabel={tiempo (min)}, ylabel={$N(t)$},
      xmin=0, xmax=5000, ymin=0, ymax=520,
      legend pos=north west]
      \addplot[thick, mobiblue] coordinates {
        (0,0) (200,5) (500,18) (1000,52) (1500,103) (2000,170)
        (2500,238) (3000,305) (3500,372) (4000,432) (4500,484) (5000,520)
      };
      \addlegendentry{$\rho=0{,}998$ — divergencia}
      \addplot[thick, deepgreen, dashed] coordinates {
        (0,0) (200,1) (500,1) (1000,1) (1500,1) (2000,1) (2500,1)
        (3000,1) (3500,1) (4000,1) (4500,1) (5000,1)
      };
      \addlegendentry{$\rho=0{,}60$ — estable}
    \end{axis}
  \end{tikzpicture}

  \vspace{2pt}
  \begin{resultadobox}[Verificación cerrada]
    \footnotesize
    Todas las pruebas pasan. El código implementa fielmente el modelo conceptual
    $M/M/1$. Procedemos a \emph{validar} la salida.
  \end{resultadobox}
\end{frame}

% ------------------------------------------------------------
\begin{frame}{Cierre Sesión 1}
  \begin{block}{Lo que hemos hecho}
    Hemos cubierto las tres familias de técnicas formales para verificar el
    modelo computacional:
  \end{block}

  \vspace{6pt}
  \begin{enumerate}\itemsep2pt
    \item \textbf{\emph{Debugging}} con aserciones de invariantes y
      \emph{logging} estructurado.
    \item \textbf{Trazas} con verificación contra cálculo manual y chequeo
      automático de invariantes.
    \item \textbf{Pruebas unitarias} con patrón AAA, incluyendo pruebas
      contra resultados analíticos.
  \end{enumerate}

  \vspace{6pt}
  \begin{alertblock}{Próxima sesión}
    Suponer ahora que el modelo está \emph{verificado} (el código hace lo que
    dice el modelo conceptual).\\
    \textbf{¿Pero el modelo conceptual representa al sistema real?}\\
    Esa es la pregunta de la \textbf{validación}.
  \end{alertblock}
\end{frame}

% =============================================================
\section*{}
\begin{frame}[plain,c]
  \begin{center}
    {\Huge \color{princetonorange}\textbf{Sesión 2}}\\[0.5cm]
    {\Large Calibración y Validación}\\
    {\large del modelo conceptual}\\[0.4cm]
    {\normalsize \emph{Face validity} $\,\cdot\,$ Validación de supuestos
      $\,\cdot\,$ Calibración\\
      Input--Output $\,\cdot\,$ Test de Turing $\,\cdot\,$ Sensibilidad}\\[0.4cm]
    {\color{mobigold}\rule{0.5\textwidth}{1.2pt}}
  \end{center}
\end{frame}

% =============================================================
\section{Sesión 2: Validación}
% =============================================================

\begin{frame}{Niveles de validación (Sargent, 2013)}
  \begin{center}
    \begin{tikzpicture}[node distance=0.4cm]
      \node[draw=princetonorange, very thick, rounded corners,
            fill=mobilightgold, minimum width=11.5cm, minimum height=0.9cm,
            font=\small] (op)
        {\textbf{Validación operacional / Input--Output}: salidas del modelo
          $\approx$ salidas reales};
      \node[draw=princetonorange, very thick, rounded corners,
            fill=mobilightgold!80, minimum width=11.5cm, minimum height=0.9cm,
            font=\small, below=of op] (con)
        {\textbf{Validación conceptual}: estructura y supuestos del modelo
          $\approx$ realidad};
      \node[draw=princetonorange, very thick, rounded corners,
            fill=mobilightgold!60, minimum width=11.5cm, minimum height=0.9cm,
            font=\small, below=of con] (datos)
        {\textbf{Validación de datos}: distribuciones de \emph{inputs}
          ajustadas correctamente};
      \node[draw=princetonorange, very thick, rounded corners,
            fill=mobilightgold!40, minimum width=11.5cm, minimum height=0.9cm,
            font=\small, below=of datos] (face)
        {\textbf{\emph{Face validity}}: el modelo es razonable a juicio de expertos};
    \end{tikzpicture}
  \end{center}

  \vspace{6pt}
  \begin{notabox}[Filosofía de fondo]
    \footnotesize
    Ningún test individual ``valida'' el modelo. La validación es una
    \textbf{acumulación de evidencia} (Balci, 1997). Cada técnica aporta un
    argumento. La confianza en el modelo crece con la cantidad y diversidad
    de pruebas superadas.
  \end{notabox}
\end{frame}

% =============================================================
\subsection{Face Validity}
% =============================================================

\begin{frame}{\emph{Face Validity}: protocolo formal}
  \begin{definicionbox}[\emph{Face validity}]
    \footnotesize
    Evaluación cualitativa del modelo por parte de \textbf{expertos del dominio}
    (no necesariamente modeladores) sobre si la estructura, los supuestos y el
    comportamiento del modelo \emph{aparentan} ser razonables.
  \end{definicionbox}

  \vspace{4pt}
  \begin{metodobox}[Protocolo de \emph{face validity} (Sargent)]
    \footnotesize
    \begin{enumerate}\itemsep1pt
      \item Identificar 3--5 expertos del sistema real (operadores, gerentes, ingenieros).
      \item Presentarles, por separado:
        \begin{itemize}\scriptsize\itemsep0pt
          \item Diagrama de flujo del modelo conceptual.
          \item Lista explícita de supuestos.
          \item Animación o resultados gráficos del modelo.
        \end{itemize}
      \item Aplicar un cuestionario estructurado.
      \item Documentar discrepancias y \textbf{re-modelar} antes de continuar.
    \end{enumerate}
  \end{metodobox}

  \vspace{4pt}
  \footnotesize
  \textbf{Ejemplo (call center):} el supervisor observa la animación. Detecta que
  el modelo nunca refleja ``rebalanceos de agentes a otra línea''. Esto cambia
  la \emph{estructura} del modelo, no solo sus parámetros.
\end{frame}

% ------------------------------------------------------------
\begin{frame}{Cuestionario estructurado para \emph{face validity}}
  \footnotesize
  \textbf{Plantilla} (puntuación 1 = totalmente en desacuerdo;
  5 = totalmente de acuerdo):

  \vspace{2pt}
  \begin{center}
    \renewcommand{\arraystretch}{1.2}
    \begin{tabular}{p{10cm}c}
      \toprule
      \textbf{Ítem} & \textbf{Puntaje} \\
      \midrule
      1.\ Las entidades modeladas (clientes, agentes, recursos) corresponden
          a las que existen en el sistema real.
        & \rule[-2pt]{0.6cm}{0.4pt} \\
      2.\ Los eventos modelados cubren los procesos relevantes del sistema.
        & \rule[-2pt]{0.6cm}{0.4pt} \\
      3.\ La política de operación implementada coincide con la real.
        & \rule[-2pt]{0.6cm}{0.4pt} \\
      4.\ Los volúmenes y tiempos observados en la animación parecen realistas.
        & \rule[-2pt]{0.6cm}{0.4pt} \\
      5.\ Los casos extremos (picos, fallas, ausencias) están razonablemente
          representados.
        & \rule[-2pt]{0.6cm}{0.4pt} \\
      6.\ Las exclusiones del modelo (qué \emph{no} se modela) están justificadas.
        & \rule[-2pt]{0.6cm}{0.4pt} \\
      \bottomrule
    \end{tabular}
  \end{center}

  \vspace{4pt}
  \begin{notabox}[Criterio]
    \footnotesize
    Promedio $\geq 4{,}0$ \textbf{y} ningún ítem por debajo de 3 (todos los
    expertos). De lo contrario, revisar el modelo.\\
    \emph{Inspirado en Sargent (2013) y Balci (1997).}
  \end{notabox}
\end{frame}

% =============================================================
\subsection{Validación de supuestos}
% =============================================================

\begin{frame}{Validación de supuestos: estructurales y de datos}
  \footnotesize
  \textbf{Dos categorías formales} (Law, cap.\ 5):

  \vspace{2pt}
  \begin{columns}[T]
    \begin{column}{0.48\textwidth}
      \begin{block}{Supuestos estructurales}
        \footnotesize
        Sobre la \emph{lógica} del sistema:
        \begin{itemize}\itemsep0pt
          \item Disciplina de cola (FIFO, prioridad).
          \item Independencia entre recursos.
          \item Bloqueo, abandonos, reintentos.
          \item Política de despacho.
        \end{itemize}
        \textbf{Validación:} observación directa, entrevistas con operadores,
        manuales de operación.
      \end{block}
    \end{column}
    \begin{column}{0.48\textwidth}
      \begin{alertblock}{Supuestos sobre datos}
        \footnotesize
        Sobre las \emph{distribuciones} de los \emph{inputs}:
        \begin{itemize}\itemsep0pt
          \item Tiempos entre llegadas $\sim\mathrm{Exp}(\lambda)$.
          \item Tiempos de servicio $\sim\mathrm{LogN}(\mu,\sigma)$.
          \item Independencia entre observaciones.
          \item Estacionariedad.
        \end{itemize}
        \textbf{Validación:} pruebas estadísticas formales
        (\emph{goodness-of-fit}).
      \end{alertblock}
    \end{column}
  \end{columns}

  \vspace{4pt}
  \begin{notabox}[Lazo con T3.1]
    \footnotesize
    Las pruebas $\chi^2$, Kolmogorov--Smirnov y Anderson--Darling vistas en T3.1
    se usan ahora con un \textbf{propósito explícito de validación}: aceptar o
    rechazar la hipótesis estructural sobre los datos del sistema real.
  \end{notabox}
\end{frame}

% ------------------------------------------------------------
\begin{frame}{Validación estructural — sucursal bancaria}
  \footnotesize
  \textbf{Sistema:} sucursal con 3 cajeros, una sola cola común.

  \vspace{2pt}
  \textbf{Supuestos del modelo conceptual:}
  \begin{enumerate}\itemsep0pt
    \item Cola única, FIFO.
    \item Cuando un cajero queda libre, toma al primero de la cola.
    \item Servicios independientes entre cajeros.
    \item No hay abandonos.
  \end{enumerate}

  \vspace{4pt}
  \textbf{Procedimiento de validación estructural:}
  \begin{center}
    \renewcommand{\arraystretch}{1.2}
    \begin{tabular}{p{4.5cm} p{8.5cm}}
      \toprule
      \textbf{Supuesto} & \textbf{Cómo validar} \\
      \midrule
      Cola única FIFO          & Observar 2 h en sucursal: ¿hay corte de fila?
                                 ¿saltos? \\
      Cajero libre toma al primero
                               & Observar transiciones: ¿algún cajero llama por
                                 número? \\
      Servicios independientes & Comparar tiempos de servicio entre cajeros
                                 (ANOVA). \\
      Sin abandonos            & Conteo de clientes que entran vs.\ que salen
                                 sin servicio. \\
      \bottomrule
    \end{tabular}
  \end{center}

  \vspace{4pt}
  \begin{alertabox}[Hallazgo típico]
    \footnotesize
    Se detecta un 4\% de abandonos $\Rightarrow$ \textbf{el supuesto 4 falla},
    hay que modelar abandonos. Esto no se descubre con pruebas estadísticas;
    se descubre \emph{mirando el sistema}.
  \end{alertabox}
\end{frame}

% ------------------------------------------------------------
\begin{frame}{Validación de datos: prueba $\chi^2$ de bondad de ajuste}
  \begin{definicionbox}[Prueba $\chi^2$]
    \footnotesize
    \textbf{Hipótesis:} $H_0$: los datos provienen de la distribución $F_0$
    propuesta.\\
    \textbf{Estadístico:}
    \[
      \chi^2_0 = \sum_{i=1}^{k} \frac{(O_i - E_i)^2}{E_i}
    \]
    $O_i$ = frecuencia observada en la clase $i$;
    $E_i = n\,p_i$ = frecuencia esperada bajo $F_0$.\\
    \textbf{Distribución:} $\chi^2_0 \sim \chi^2_{k-s-1}$ bajo $H_0$, con
    $s$ = parámetros estimados de los datos.\\
    \textbf{Decisión:} rechazar $H_0$ si $\chi^2_0 > \chi^2_{\alpha,\,k-s-1}$.
  \end{definicionbox}

  \vspace{4pt}
  \begin{notabox}[Reglas prácticas]
    \footnotesize
    (i) cada $E_i \geq 5$;\quad
    (ii) usar al menos $k=5$ clases;\quad
    (iii) preferir clases de \emph{probabilidad igual} bajo $F_0$.
  \end{notabox}
\end{frame}

% ------------------------------------------------------------
\begin{frame}{$\chi^2$ trabajado — ¿interarribos $\sim$ Exponencial?}
  \footnotesize
  \textbf{Datos:} $n=100$ tiempos entre llegadas a un cajero. Media muestral
  $\bar X=2{,}50$ min $\Rightarrow \hat\lambda=1/\bar X=0{,}40$.

  \vspace{2pt}
  \begin{columns}[T]
    \begin{column}{0.48\textwidth}
      \textbf{Construcción de clases con probabilidades iguales}
      ($k=5$, $p_i=0{,}20$):
      \begin{center}\scriptsize
        \renewcommand{\arraystretch}{1.15}
        \begin{tabular}{cccccc}
          \toprule
          $i$ & Intervalo & $p_i$ & $E_i$ & $O_i$ & $\frac{(O-E)^2}{E}$ \\
          \midrule
          1 & $[0;\;0{,}558)$        & 0.20 & 20 & 23 & 0.450 \\
          2 & $[0{,}558;\;1{,}277)$  & 0.20 & 20 & 17 & 0.450 \\
          3 & $[1{,}277;\;2{,}291)$  & 0.20 & 20 & 18 & 0.200 \\
          4 & $[2{,}291;\;4{,}023)$  & 0.20 & 20 & 24 & 0.800 \\
          5 & $[4{,}023;\;\infty)$   & 0.20 & 20 & 18 & 0.200 \\
          \midrule
          \multicolumn{5}{r}{$\chi^2_0=$} & \textbf{2.100} \\
          \bottomrule
        \end{tabular}
      \end{center}

      \vspace{2pt}
      Con $\alpha=0{,}05$, $k=5$, $s=1$: $\chi^2_{0{,}05,\,3}=7{,}815$.
      Como $2{,}10 < 7{,}815$ \textbf{no se rechaza $H_0$}.
    \end{column}
    \begin{column}{0.48\textwidth}
      \begin{tikzpicture}
        \begin{axis}[width=\textwidth, height=4.6cm,
          title={Densidad $\chi^2_3$ y región de rechazo},
          xlabel={$x$}, ylabel={$f(x)$},
          xmin=0, xmax=14, ymin=0, ymax=0.27,
          legend pos=north east]
          \addplot[name path=A, thick, mobiblue, domain=0.01:14, samples=160]
            {(1/(2^1.5*1.7725))*x^0.5*exp(-x/2)};
          \addlegendentry{$\chi^2_3$}
          \addplot[name path=B, draw=none, domain=7.815:14, samples=80]
            {(1/(2^1.5*1.7725))*x^0.5*exp(-x/2)};
          \addplot[fill=princetonorange!30] fill between[of=A and B,
            soft clip={domain=7.815:14}];
          \draw[thick, princetonorange, dashed] (axis cs:7.815,0)
            -- (axis cs:7.815,0.05) node[above, font=\tiny]{$\chi^2_{0{,}05}$};
          \draw[thick, mobiblue] (axis cs:2.10,0)
            -- (axis cs:2.10,0.18) node[above, font=\tiny]{$\chi^2_0=2{,}10$};
        \end{axis}
      \end{tikzpicture}
      \vspace{-2pt}
      {\centering\scriptsize El estadístico cae lejos de la región crítica.\par}
    \end{column}
  \end{columns}

  \vspace{2pt}
  \begin{resultadobox}[Conclusión]
    \footnotesize
    \textcolor{deepgreen}{El supuesto exponencial es razonable para los interarribos}
    a este nivel de evidencia.
  \end{resultadobox}
\end{frame}

% ------------------------------------------------------------
\begin{frame}{Validación de datos: prueba de Kolmogorov--Smirnov}
  \begin{definicionbox}[Prueba KS]
    \footnotesize
    \textbf{Estadístico:}
    \[
      D_n = \sup_x \big| F_n(x) - F_0(x) \big|
    \]
    $F_n$ = CDF empírica; $F_0$ = CDF teórica propuesta.\\
    \textbf{Decisión:} rechazar $H_0$ si $D_n > D_{\alpha,n}$ (valor crítico).
  \end{definicionbox}

  \vspace{4pt}
  \footnotesize
  \textbf{Ventajas sobre $\chi^2$:}
  \begin{itemize}\itemsep1pt
    \item No requiere agrupación en clases (válido para muestras pequeñas).
    \item Más potente para distribuciones continuas.
  \end{itemize}

  \vspace{4pt}
  \begin{alertabox}[Limitación]
    \footnotesize
    Cuando los parámetros se estiman de los datos (Exp, Normal, Weibull) los
    valores críticos tabulados son \textbf{conservadores}. Hay tablas corregidas:
    \textbf{Lilliefors} (normal), \textbf{KS para exponencial con} $\hat\lambda$.
  \end{alertabox}
\end{frame}

% ------------------------------------------------------------
\begin{frame}{KS trabajado — muestra pequeña de tiempos de servicio}
  \footnotesize
  \textbf{Datos:} $n=10$ tiempos de servicio (min):
  $\{0{,}45;\;0{,}80;\;1{,}20;\;1{,}55;\;2{,}10;\;2{,}40;\;3{,}05;\;3{,}80;\;4{,}50;\;6{,}15\}$.\\
  $\bar X=2{,}60 \Rightarrow \hat\lambda=0{,}3846$. $H_0$: Exp$(\hat\lambda)$.

  \vspace{2pt}
  \begin{columns}[T]
    \begin{column}{0.48\textwidth}
      \begin{center}\scriptsize
        \renewcommand{\arraystretch}{1.15}
        \begin{tabular}{cccccc}
          \toprule
          $i$ & $x_{(i)}$ & $F_n$ & $F_0$ & $|F_n-F_0|$ & $|F_n^- - F_0|$ \\
          \midrule
          1  & 0.45 & 0.10 & 0.1585 & 0.0585 & 0.1585 \\
          2  & 0.80 & 0.20 & 0.2659 & 0.0659 & 0.1659 \\
          3  & 1.20 & 0.30 & 0.3700 & 0.0700 & 0.1700 \\
          4  & 1.55 & 0.40 & 0.4493 & 0.0493 & 0.1493 \\
          5  & 2.10 & 0.50 & 0.5538 & 0.0538 & 0.1538 \\
          6  & 2.40 & 0.60 & 0.6020 & 0.0020 & 0.1020 \\
          7  & 3.05 & 0.70 & 0.6904 & 0.0096 & 0.0904 \\
          8  & 3.80 & 0.80 & 0.7676 & 0.0324 & 0.0676 \\
          9  & 4.50 & 0.90 & 0.8254 & 0.0746 & 0.0254 \\
          10 & 6.15 & 1.00 & 0.9062 & 0.0938 & 0.0062 \\
          \bottomrule
        \end{tabular}
      \end{center}

      \vspace{2pt}
      $D_{10}=0{,}1700$. Crítico (Lilliefors-exp): $D_{0{,}05;\,10}\approx 0{,}338$.\\
      $0{,}170 < 0{,}338 \Rightarrow$
      \textcolor{deepgreen}{no se rechaza $H_0$}.
    \end{column}
    \begin{column}{0.48\textwidth}
      \begin{tikzpicture}
        \begin{axis}[width=\textwidth, height=4.6cm,
          title={ECDF $F_n$ vs.\ teórica $F_0$},
          xlabel={$x$ (min)}, ylabel={$F(x)$},
          xmin=0, xmax=7, ymin=0, ymax=1.05,
          legend pos=south east]
          \addplot[const plot, thick, mobiblue, mark=none] coordinates {
            (0,0) (0.45,0) (0.45,0.10) (0.80,0.10) (0.80,0.20)
            (1.20,0.20) (1.20,0.30) (1.55,0.30) (1.55,0.40)
            (2.10,0.40) (2.10,0.50) (2.40,0.50) (2.40,0.60)
            (3.05,0.60) (3.05,0.70) (3.80,0.70) (3.80,0.80)
            (4.50,0.80) (4.50,0.90) (6.15,0.90) (6.15,1.00) (7,1.00)
          };
          \addlegendentry{$F_n$}
          \addplot[thick, princetonorange, domain=0:7, samples=120]
            {1-exp(-0.3846*x)};
          \addlegendentry{Exp$(0{,}3846)$}
          \draw[thick, deepgreen, dashed] (axis cs:1.20,0.30)
            -- (axis cs:1.20,0.47);
          \node[font=\tiny, deepgreen, right] at (axis cs:1.25,0.38) {$D_{10}$};
        \end{axis}
      \end{tikzpicture}
    \end{column}
  \end{columns}
\end{frame}

% =============================================================
\subsection{Calibración}
% =============================================================

\begin{frame}{Calibración: definición y proceso}
  \begin{definicionbox}[Calibración]
    \footnotesize
    Ajuste iterativo de los \textbf{parámetros} y, ocasionalmente, de la
    \textbf{estructura} del modelo, hasta que sus salidas reproducen
    razonablemente las del sistema real bajo condiciones conocidas.
  \end{definicionbox}

  \vspace{4pt}
  \footnotesize
  \textbf{Calibración $\neq$ Validación:}
  \begin{itemize}\itemsep1pt
    \item \textbf{Calibrar:} \emph{ajustar} para que el modelo se parezca a los datos.
    \item \textbf{Validar:} \emph{probar} el modelo contra datos
      \textbf{independientes} a los usados para calibrar.
  \end{itemize}

  \vspace{6pt}
  \begin{metodobox}[Procedimiento de división de muestra]
    \footnotesize
    \begin{enumerate}\itemsep1pt
      \item Particionar los datos en \textbf{conjunto de calibración} (70\%) y
        \textbf{conjunto de validación} (30\%).
      \item Ajustar parámetros sobre el conjunto de calibración (MLE).
      \item Correr la simulación y comparar con el conjunto de validación.
      \item Si el desempeño en validación es pobre, revisar
        \emph{estructura} (no solo parámetros) y repetir.
    \end{enumerate}
  \end{metodobox}
\end{frame}

% ------------------------------------------------------------
\begin{frame}{Ejemplo de calibración — tiempos de servicio del cajero}
  \footnotesize
  \textbf{Sistema:} sucursal bancaria. $n=600$ observaciones de tiempos de servicio.

  \vspace{2pt}
  \textbf{Paso 1.\ Partición:}\\
  Calibración: 420 obs (70\%).\quad Validación: 180 obs (30\%).

  \vspace{4pt}
  \textbf{Paso 2.\ Calibración por máxima verosimilitud:}\\
  Suponemos $\mathrm{LogN}(\mu,\sigma)$. MLE sobre el conjunto de calibración:\\
  $\hat\mu=0{,}812$, $\hat\sigma=0{,}443$ \quad ($\bar X_{\text{cal}}=2{,}48$ min).

  \vspace{4pt}
  \textbf{Paso 3.\ Validación cruzada:} prueba KS de los 180 datos de validación
  contra $\mathrm{LogN}(0{,}812;\,0{,}443)$.\\
  $D_{180}=0{,}061$. Crítico Lilliefors $=0{,}066$ a $\alpha=0{,}05$.\\
  \textcolor{deepgreen}{$0{,}061 < 0{,}066 \Rightarrow$ ajuste valida en datos no usados.}

  \vspace{6pt}
  \begin{alertabox}[Advertencia]
    \footnotesize
    Si la prueba KS hubiera fallado en la \emph{misma} muestra usada para calibrar,
    no significaría nada (autocomplacencia). La prueba en \textbf{datos
    independientes} es lo que da validez.
  \end{alertabox}
\end{frame}

% =============================================================
\subsection{Validación Input-Output}
% =============================================================

\begin{frame}{Validación Input--Output: idea general}
  \begin{definicionbox}[Validación I/O]
    \footnotesize
    Comparar las \textbf{salidas} del modelo de simulación con las
    \textbf{salidas observadas} del sistema real, cuando ambos operan bajo
    \textbf{las mismas condiciones de entrada}.
  \end{definicionbox}

  \vspace{6pt}
  \begin{center}
    \begin{tikzpicture}[node distance=0.8cm, >=Stealth]
      \node[draw=mobiblue, fill=mobilightblue, rounded corners,
            minimum width=2.5cm, minimum height=1cm, font=\footnotesize] (in)
        {Mismos \emph{inputs}};
      \node[draw=mobiblue, fill=mobilightblue, rounded corners,
            minimum width=3cm, minimum height=1cm, right=of in,
            yshift=1cm, font=\footnotesize] (real) {Sistema real};
      \node[draw=mobiblue, fill=mobilightblue, rounded corners,
            minimum width=3cm, minimum height=1cm, right=of in,
            yshift=-1cm, font=\footnotesize] (sim) {Simulación};
      \node[draw=princetonorange, fill=mobilightgold, rounded corners,
            minimum width=3.5cm, minimum height=1.2cm, align=center,
            right=of real, yshift=-1cm, xshift=0.5cm, font=\footnotesize]
            (test) {\textbf{Prueba estadística}\\¿$Y_{\text{real}}\approx Y_{\text{sim}}$?};
      \draw[->, thick, mobiblue] (in) -- (real);
      \draw[->, thick, mobiblue] (in) -- (sim);
      \draw[->, thick, mobiblue] (real) --
        node[above, font=\scriptsize]{$Y_{\text{real}}$}
        ([yshift=2pt]test.west|-real);
      \draw[->, thick, mobiblue] (sim) --
        node[below, font=\scriptsize]{$Y_{\text{sim}}$}
        ([yshift=-2pt]test.west|-sim);
    \end{tikzpicture}
  \end{center}

  \vspace{4pt}
  \footnotesize\textbf{Pruebas estadísticas comunes:}
  \begin{itemize}\itemsep1pt
    \item Prueba $t$ \textbf{pareada} (cuando hay correspondencia uno-a-uno
      entre días/horas).
    \item Prueba $t$ de \textbf{dos muestras} (cuando no la hay).
    \item IC para la diferencia de medias: si $0\in IC \Rightarrow$ no hay
      evidencia de diferencia.
  \end{itemize}
\end{frame}

% ------------------------------------------------------------
\begin{frame}{I/O con prueba $t$ pareada — definición formal}
  \begin{formulabox}[Prueba $t$ pareada]
    \footnotesize
    Sean $Y_i^{\text{real}}, Y_i^{\text{sim}}$ con $i=1,\ldots,n$ pares
    (mismo día/hora). Definir
    $D_i = Y_i^{\text{real}}-Y_i^{\text{sim}}$, $\bar D$ y $S_D$ media y
    desviación muestrales.
    \[
      T = \frac{\bar D}{S_D/\sqrt{n}} \sim t_{n-1}\;\;\text{ bajo }\;
      H_0:\mu_D=0
    \]
    IC del $100(1-\alpha)\%$ para $\mu_D$:
    \[
      \bar D \pm t_{\alpha/2,\,n-1}\,\frac{S_D}{\sqrt{n}}
    \]
  \end{formulabox}

  \vspace{6pt}
  \begin{notabox}[Por qué pareada]
    \footnotesize
    Controla la variabilidad común a real y simulación (mismo día = mismas
    condiciones), aumentando la potencia frente a la prueba de dos muestras.
  \end{notabox}
\end{frame}

% ------------------------------------------------------------
\begin{frame}{I/O trabajado — tiempo medio de espera diario, $n=10$ días}
  \footnotesize
  \textbf{Sistema:} cola de servicio al cliente. Para cada día se mide el tiempo
  medio de espera real y se ejecuta una réplica con los \emph{inputs} de ese día.

  \vspace{2pt}
  \begin{center}\scriptsize
    \renewcommand{\arraystretch}{1.2}
    \begin{tabular}{ccccccccccc}
      \toprule
      Día $i$ & 1 & 2 & 3 & 4 & 5 & 6 & 7 & 8 & 9 & 10 \\
      \midrule
      $Y_i^{\text{real}}$ (min) & 4.2 & 5.1 & 3.8 & 6.0 & 4.7 & 5.5 & 3.9 & 4.4 & 5.8 & 4.6 \\
      $Y_i^{\text{sim}}$  (min) & 4.5 & 4.8 & 3.6 & 5.7 & 4.9 & 5.2 & 4.1 & 4.2 & 5.5 & 4.8 \\
      $D_i$                     & -0.3 & 0.3 & 0.2 & 0.3 & -0.2 & 0.3 & -0.2 & 0.2 & 0.3 & -0.2 \\
      \bottomrule
    \end{tabular}
  \end{center}

  \vspace{4pt}
  \begin{columns}[T]
    \begin{column}{0.48\textwidth}
      \footnotesize
      \textbf{Cálculos:}
      \[
        \bar D=0{,}07,\;\; S_D=0{,}2627,\;\; n=10
      \]
      \[
        T=\frac{0{,}07}{0{,}2627/\sqrt{10}}=\frac{0{,}07}{0{,}0831}=0{,}843
      \]
      Con $\alpha=0{,}05$, $t_{0{,}025;9}=2{,}262$:
      $|T|=0{,}843 < 2{,}262$ $\Rightarrow$ no se rechaza $H_0$.\\[2pt]
      IC 95\% para $\mu_D$:
      $0{,}07 \pm 2{,}262\cdot 0{,}0831 = [-0{,}118;\;0{,}258]$.\\
      Como $0\in IC$,
      \textcolor{deepgreen}{\textbf{no hay evidencia de diferencia}}.
    \end{column}
    \begin{column}{0.48\textwidth}
      \begin{tikzpicture}
        \begin{axis}[width=\textwidth, height=4.4cm,
          title={Diferencias diarias $D_i$ y su IC},
          xlabel={Día}, ylabel={$D_i$ (min)},
          xmin=0, xmax=11, ymin=-0.4, ymax=0.4,
          legend pos=south east]
          \addplot[only marks, mark=*, mark size=1.6pt, mobiblue]
            coordinates {
              (1,-0.3) (2,0.3) (3,0.2) (4,0.3) (5,-0.2)
              (6,0.3) (7,-0.2) (8,0.2) (9,0.3) (10,-0.2)
            };
          \addlegendentry{$D_i$}
          \addplot[thick, princetonorange, dashed, domain=0:11] {0.07};
          \addlegendentry{$\bar D$}
          \addplot[thick, mobigold, domain=0:11] {-0.118};
          \addplot[thick, mobigold, domain=0:11] {0.258};
        \end{axis}
      \end{tikzpicture}
      \vspace{-2pt}
      {\centering\scriptsize El IC contiene $0$.\par}
    \end{column}
  \end{columns}
\end{frame}

% ------------------------------------------------------------
\begin{frame}{Significancia estadística vs.\ operacional}
  \footnotesize
  \textbf{Caso opuesto:} con $n$ muy grande, casi cualquier diferencia se vuelve
  estadísticamente significativa. ¿Significa eso que el modelo no es válido?

  \vspace{4pt}
  \begin{metodobox}[Criterio adicional: tolerancia operacional $\delta^*$]
    \footnotesize
    Definir \emph{a priori} un umbral $\delta^*$ de diferencia que es
    \textbf{operacionalmente irrelevante}.
    \begin{itemize}\itemsep1pt
      \item $IC \subset (-\delta^*,\delta^*)$ $\Rightarrow$
        \textbf{validación fuerte} (diferencia indistinguible de irrelevante).
      \item $0 \in IC$ pero $IC \not\subset (-\delta^*,\delta^*)$ $\Rightarrow$
        \textbf{validación débil} (no se detecta diferencia, pero podría haberla
        relevante; aumentar $n$).
      \item $0 \notin IC$ y $IC \subset (-\delta^*,\delta^*)$ $\Rightarrow$
        \textbf{diferencia detectada pero irrelevante}.
      \item $0 \notin IC$ y $IC \not\subset (-\delta^*,\delta^*)$ $\Rightarrow$
        \textbf{revisar el modelo}.
    \end{itemize}
  \end{metodobox}

  \vspace{4pt}
  \begin{alertabox}[El criterio que separa al modelador del estudiante]
    \footnotesize
    Esta clasificación es la que separa al modelador entrenado del estudiante
    que mira solo el $p$-valor (Law, cap.\ 5).
  \end{alertabox}
\end{frame}

% =============================================================
\subsection{Test de Turing}
% =============================================================

\begin{frame}{Test de Turing para simulación (Schruben, 1980)}
  \begin{definicionbox}[Test de Turing]
    \footnotesize
    Se le presentan a un experto del sistema $K$ reportes mezclados, algunos
    generados por el sistema real y otros por la simulación. El experto debe
    clasificar cada reporte como ``real'' o ``simulado''. Si el experto
    \textbf{no puede distinguir} mejor que el azar, el modelo se considera válido.
  \end{definicionbox}

  \vspace{4pt}
  \begin{metodobox}[Procedimiento de Schruben]
    \footnotesize
    \begin{enumerate}\itemsep1pt
      \item Generar $K_R$ reportes del sistema real y $K_S$ reportes de la
        simulación, con \emph{el mismo formato}.
      \item Mezclarlos aleatoriamente; etiquetar internamente.
      \item Presentar al experto, sin decirle cuáles son cuáles.
      \item El experto clasifica cada reporte.
      \item Construir tabla de aciertos $\to$ prueba estadística.
    \end{enumerate}
  \end{metodobox}
\end{frame}

% ------------------------------------------------------------
\begin{frame}{Test de Turing: análisis estadístico de los aciertos}
  \footnotesize
  \textbf{Diseño $K_R=K_S=K$ con $K=10$ reportes de cada tipo (total 20).}

  \vspace{2pt}
  \begin{columns}[T]
    \begin{column}{0.48\textwidth}
      \textbf{Tabla de contingencia} (clasificación del experto):
      \begin{center}\scriptsize
        \renewcommand{\arraystretch}{1.2}
        \begin{tabular}{l|cc|c}
          \toprule
          & Clasif.\ ``Real'' & Clasif.\ ``Sim.'' & Total \\
          \midrule
          Realmente real      & 6 & 4 & 10 \\
          Realmente simulado  & 5 & 5 & 10 \\
          \midrule
          Total               & 11 & 9 & 20 \\
          \bottomrule
        \end{tabular}
      \end{center}

      \vspace{2pt}
      \textbf{Aciertos:} $6+5=11$ de 20 $\Rightarrow$ tasa $=0{,}55$.

      \vspace{4pt}
      \textbf{Prueba binomial:} $H_0$: $p=0{,}5$ (el experto adivina).
      \[
        \Prob(X\geq 11\mid n=20,\,p=0{,}5) = 0{,}412
      \]
      $p$-valor alto $\Rightarrow$
      \textcolor{deepgreen}{no se rechaza $H_0$}: el experto no logra distinguir
      mejor que el azar.
    \end{column}
    \begin{column}{0.48\textwidth}
      \begin{tikzpicture}
        \begin{axis}[width=\textwidth, height=4.6cm,
          title={Distribución de aciertos bajo $H_0$ ($n=20$, $p=0{,}5$)},
          xlabel={aciertos $X$}, ylabel={$\Prob(X=k)$},
          xmin=4, xmax=16, ymin=0, ymax=0.20,
          ybar, bar width=6pt, legend pos=north east]
          \addplot[fill=mobiblue!50, draw=mobiblue] coordinates {
            (5,0.0148) (6,0.0370) (7,0.0739) (8,0.1201) (9,0.1602)
            (10,0.1762) (11,0.1602) (12,0.1201) (13,0.0739) (14,0.0370)
            (15,0.0148)
          };
          \addlegendentry{Binomial$(20,0{,}5)$}
          \addplot[fill=princetonorange!50, draw=princetonorange]
            coordinates {(11,0.1602)};
          \addlegendentry{$X=11$ obs.}
        \end{axis}
      \end{tikzpicture}
    \end{column}
  \end{columns}

  \vspace{4pt}
  \begin{resultadobox}[Conclusión]
    \footnotesize
    Las salidas simuladas son \textbf{indistinguibles}, en juicio experto, de
    las reales.
  \end{resultadobox}
\end{frame}

% ------------------------------------------------------------
\begin{frame}{Buenas prácticas del Test de Turing}
  \footnotesize
  \begin{itemize}\itemsep2pt
    \item Usar formato \textbf{idéntico} en reportes reales y simulados (incluso
      números de decimales y orden de filas). Cualquier diferencia de formato
      es una pista que invalida el test.
    \item Usar \textbf{varios expertos independientes}, no uno solo.
    \item Pedirle al experto que \textbf{justifique} su clasificación: las
      razones revelan qué aspectos del modelo son débiles, incluso cuando el
      test pasa.
    \item Repetir el test con $K$ moderadamente grande (al menos 10 de cada
      tipo) para tener potencia estadística.
    \item No es un test único: combinar con I/O y validación de supuestos.
  \end{itemize}

  \vspace{6pt}
  \begin{alertabox}[Cuidado]
    \footnotesize
    Un test de Turing aprobado \textbf{no} significa que el modelo sea correcto,
    solo que es \emph{indistinguible} en los \emph{outputs} examinados. El
    experto puede pasar por alto patrones sutiles.
  \end{alertabox}
\end{frame}

% =============================================================
% EJEMPLO 2 — Validación que detecta un bug sutil
% =============================================================
\section{Ejemplo 2 (intermedio)}

\begin{frame}{Ejemplo 2 — Validación del $M/M/1$ y un bug sutil}
  \begin{ejemplobox}[Sistema y dos versiones del código]
    \footnotesize
    Mismo sistema $M/M/1$ con $\lambda=3$, $\mu=5$ que el Ejemplo 1.
    Comparamos:
    \begin{itemize}\itemsep1pt
      \item \textbf{Versión correcta:} $\mu=5$ en el código.
      \item \textbf{Versión con bug:} el programador escribió $\mu=4$ por
        error de transcripción.
    \end{itemize}
    Ambas pasan la verificación estructural sin problemas.
  \end{ejemplobox}

  \vspace{4pt}
  \footnotesize\textbf{Resultados de $n_R=30$ réplicas con $T=2000$ min:}
  \begin{center}\scriptsize
    \renewcommand{\arraystretch}{1.25}
    \begin{tabular}{lcccc}
      \toprule
      \textbf{Métrica} & \textbf{Analítico} & $\bar Y_{30}$ & \textbf{IC 95\%}
        & \textbf{Diagnóstico} \\
      \midrule
      \multicolumn{5}{l}{\textcolor{deepgreen}{\textbf{Versión correcta}
        ($\mu=5$, $\rho=0{,}60$)}} \\
      $W_q$ (min) & $0{,}300$ & $0{,}308$ & $[0{,}293;\;0{,}323]$
        & \textcolor{deepgreen}{Validado $\checkmark$} \\
      $L_q$ (cl.) & $0{,}900$ & $0{,}924$ & $[0{,}877;\;0{,}971]$
        & \textcolor{deepgreen}{Validado $\checkmark$} \\
      $\rho$      & $0{,}600$ & $0{,}598$ & $[0{,}596;\;0{,}600]$
        & \textcolor{deepgreen}{Validado $\checkmark$} \\
      \midrule
      \multicolumn{5}{l}{\textcolor{red!70!black}{\textbf{Versión con bug}
        ($\mu=4$, $\rho_{\text{aparente}}=0{,}75$)}} \\
      $W_q$ (min) & $0{,}300$ & $0{,}762$ & $[0{,}701;\;0{,}823]$
        & \textcolor{red!70!black}{Rechazado} \\
      $\rho$      & $0{,}600$ & $0{,}749$ & $[0{,}745;\;0{,}753]$
        & \textcolor{red!70!black}{Rechazado} \\
      \bottomrule
    \end{tabular}
  \end{center}
\end{frame}

\begin{frame}{Ejemplo 2 — Lectura del bug a partir del IC}
  \begin{columns}[T]
    \begin{column}{0.50\textwidth}
      \begin{tikzpicture}
        \begin{axis}[width=\textwidth, height=5.0cm,
          title={IC al 95\% para $\rho$},
          xbar, bar width=8pt, xmin=0.55, xmax=0.78,
          symbolic y coords={Bug $\mu=4$, Correcto $\mu=5$, Referencia},
          ytick=data, nodes near coords,
          nodes near coords style={font=\tiny},
          enlarge y limits=0.4]
          \addplot[fill=princetonorange!20, draw=princetonorange,
                   error bars/.cd, x dir=both, x explicit] coordinates {
            (0.749,Bug $\mu=4$) +- (0.004,0)
          };
          \addplot[fill=mobiblue!30, draw=mobiblue,
                   error bars/.cd, x dir=both, x explicit] coordinates {
            (0.598,Correcto $\mu=5$) +- (0.002,0)
          };
          \addplot[mark=*, mark size=2pt, deepgreen, only marks] coordinates {
            (0.600,Referencia)
          };
        \end{axis}
      \end{tikzpicture}
    \end{column}
    \begin{column}{0.48\textwidth}
      \footnotesize
      \textbf{Diagnóstico mecánico:}
      \begin{itemize}\itemsep1pt
        \item $\hat\rho_{\text{bug}} \approx 0{,}749 \approx 3/4$.
        \item Si la tasa de llegada es la documentada ($\lambda=3$),
          la tasa de servicio efectiva es $\hat\mu=\lambda/\hat\rho\approx 4$.
        \item Apunta directamente al parámetro mal cargado.
      \end{itemize}

      \vspace{4pt}
      \begin{alertabox}[Verificación pasó, validación detectó]
        \footnotesize
        El código \emph{funcionaba} correctamente — solo con el parámetro equivocado.
        Trazas, conservación y pruebas degeneradas habrían pasado idénticas.
        Sin la comparación contra una referencia analítica, el bug habría llegado
        a producción.
      \end{alertabox}
    \end{column}
  \end{columns}

  \vspace{4pt}
  \begin{resultadobox}[Resultado del Ejemplo 2]
    \footnotesize
    Errores en \textbf{parámetros} no se detectan por verificación (todo pasa);
    se detectan por \textbf{validación} (los IC no contienen la referencia).
    Por eso ambos procesos son obligatorios.
  \end{resultadobox}
\end{frame}

% =============================================================
% EJEMPLO 3 — APLICADO M/E_2/1 + SENSIBILIDAD
% =============================================================
\section{Ejemplo 3 (aplicado)}

\begin{frame}{Ejemplo 3 — Celda de manufactura $M/E_2/1$}
  \begin{columns}[T]
    \begin{column}{0.48\textwidth}
      \begin{situacionbox}[Sistema]
        \footnotesize
        Una celda de manufactura recibe piezas en proceso de Poisson con
        $\lambda=8$ pzas/h. El servicio tiene dos etapas (carga y mecanizado)
        exponenciales con $\mu_{\text{fase}}=10$ op/h cada una:
        \[
          S \sim E_2(10),\;\; \E[S]=0{,}2 \text{ h},\;\; \mathrm{CV}_S=1/\sqrt{2}
        \]
      \end{situacionbox}

      \vspace{4pt}
      \textbf{Parámetros derivados:}
      \[
        \rho=8\times 0{,}2=0{,}80
      \]
      \[
        \E[S^2]=\Var[S]+\E[S]^2=0{,}02+0{,}04=0{,}06
      \]
    \end{column}
    \begin{column}{0.48\textwidth}
      \begin{tikzpicture}[>=Stealth, node distance=1.4cm,
        proc/.style={rectangle, rounded corners=3pt, draw=mobiblue, thick,
          fill=mobilightblue, minimum width=1.4cm, minimum height=0.9cm,
          align=center, font=\scriptsize}]
        \node[proc] (q) {Cola};
        \node[proc, right=of q, fill=mobilightgold] (s1) {$F_1$\\$\mu_f=10$};
        \node[proc, right=of s1, fill=mobilightgold] (s2) {$F_2$\\$\mu_f=10$};
        \node[font=\scriptsize, left=0.6cm of q] (in) {$\lambda=8$};
        \node[font=\scriptsize, right=0.6cm of s2] (out) {salidas};
        \draw[->, thick, mobiblue] (in) -- (q);
        \draw[->, thick, mobiblue] (q) -- (s1);
        \draw[->, thick, princetonorange] (s1) -- (s2);
        \draw[->, thick, deepgreen] (s2) -- (out);
        \node[font=\tiny, below=0.0cm of s1] {Erlang-2 servidor único};
      \end{tikzpicture}

      \vspace{6pt}
      \textbf{Modelo conceptual:} cola $M/E_2/1$ (Poisson + Erlang-2 + 1 servidor).
      Aplicaremos los \textbf{6 pasos} del proceso integrado de V\&V.
    \end{column}
  \end{columns}
\end{frame}

\begin{frame}{Ejemplo 3 — Verificación}
  \begin{center}\footnotesize
    \renewcommand{\arraystretch}{1.25}
    \begin{tabular}{p{4.0cm}p{5.0cm}p{4.4cm}}
      \toprule
      \textbf{Prueba} & \textbf{Esperado} & \textbf{Resultado} \\
      \midrule
      Tráfico extremo $\lambda=9{,}9$    & $L_q\to\infty$ en $T=500$ h
                                         & $L_q>200$ pzas $\checkmark$ \\
      Simplificación: $E_2$ vs.\ $S_1+S_2$
                                         & $\E[S]=0{,}2$, $\Var=0{,}02$,
                                           asimetría $\sqrt{2}\approx 1{,}41$
                                         & Coincide al 4to decimal $\checkmark$ \\
      Conservación, $T=2000$ h           & $\approx 16\,000$ llegadas
                                         & $16\,023 = 16\,023 + 0$ $\checkmark$ \\
      Aserciones                         & Ninguna falla
                                         & 0 violaciones en $10^7$ eventos $\checkmark$ \\
      \bottomrule
    \end{tabular}
  \end{center}

  \vspace{6pt}
  \begin{tikzpicture}
    \begin{axis}[width=12cm, height=3.6cm,
      title={Distribución empírica del servicio $S$ vs.\ Erlang-2},
      xlabel={$s$ (h)}, ylabel={densidad},
      xmin=0, xmax=0.8, ymin=0, ymax=4.0,
      legend pos=north east]
      \addplot[ybar, fill=mobiblue!30, draw=mobiblue, bar width=2pt]
        coordinates {
          (0.05,2.20) (0.10,3.30) (0.15,3.50) (0.20,3.10) (0.25,2.50)
          (0.30,1.80) (0.35,1.20) (0.40,0.85) (0.45,0.55) (0.50,0.36)
          (0.55,0.22) (0.60,0.13) (0.65,0.08) (0.70,0.05) (0.75,0.03)
        };
      \addlegendentry{Histograma simulado ($n=10^4$)}
      \addplot[thick, princetonorange, domain=0:0.8, samples=120]
        {100*x*exp(-10*x)};
      \addlegendentry{Erlang-2 teórica}
    \end{axis}
  \end{tikzpicture}

  \vspace{2pt}
  \begin{notabox}[Lectura]
    \footnotesize
    El generador de tiempos de servicio implementa correctamente $S=S_1+S_2$
    con $S_i\sim\mathrm{Exp}(10)$. La verificación está cerrada.
  \end{notabox}
\end{frame}

\begin{frame}{Ejemplo 3 — Validación con Pollaczek--Khinchine}
  \begin{columns}[T]
    \begin{column}{0.48\textwidth}
      \begin{formulabox}[P-K para $M/G/1$]
        \footnotesize
        \[
          W_q^{\text{P-K}} = \frac{\lambda\,\E[S^2]}{2(1-\rho)}
        \]
      \end{formulabox}

      \vspace{2pt}
      \footnotesize
      \textbf{Cálculo:}
      \[
        W_q^{\text{P-K}} = \frac{8\times 0{,}06}{2\times 0{,}2}
                         = \frac{0{,}48}{0{,}4}
                         = 1{,}20 \text{ h}
      \]

      \vspace{4pt}
      \textbf{Comparación con $M/M/1$ equivalente} ($\E[S]=0{,}2$):
      \[
        W_q^{M/M/1} = \frac{0{,}8}{5\times 0{,}2} = 1{,}60 \text{ h}
      \]
      $\Rightarrow$ Erlang-2 reduce $W_q$ en 25\% gracias a la
      menor variabilidad del servicio.
    \end{column}
    \begin{column}{0.48\textwidth}
      \footnotesize
      \textbf{Resultados de simulación}\\
      ($n_R=50$ réplicas, $T=2000$ h)
      \begin{center}\scriptsize
        \renewcommand{\arraystretch}{1.25}
        \begin{tabular}{lccc}
          \toprule
          \textbf{Métrica} & \textbf{P-K} & \textbf{Sim.} & \textbf{IC 95\%} \\
          \midrule
          $W_q$ (h)    & $1{,}200$ & $1{,}214$ & $[1{,}176;\,1{,}252]$ \\
          $L_q$ (pzas) & $9{,}600$ & $9{,}712$ & $[9{,}408;\,10{,}016]$ \\
          $\rho$       & $0{,}800$ & $0{,}799$ & $[0{,}795;\,0{,}803]$ \\
          \bottomrule
        \end{tabular}
      \end{center}

      \vspace{4pt}
      \begin{tikzpicture}
        \begin{axis}[width=\textwidth, height=3.6cm,
          xlabel={Métrica}, ylabel={Valor},
          ybar, bar width=8pt, ymin=0, ymax=11,
          symbolic x coords={$W_q$, $L_q$, $\rho$},
          xtick=data, nodes near coords,
          nodes near coords style={font=\tiny},
          legend pos=north east]
          \addplot[fill=mobiblue!70, draw=mobiblue]
            coordinates {($W_q$,1.20) ($L_q$,9.60) ($\rho$,0.80)};
          \addlegendentry{P-K}
          \addplot[fill=princetonorange!70, draw=princetonorange]
            coordinates {($W_q$,1.21) ($L_q$,9.71) ($\rho$,0.80)};
          \addlegendentry{Sim.}
        \end{axis}
      \end{tikzpicture}
    \end{column}
  \end{columns}

  \vspace{2pt}
  \begin{resultadobox}[Validación operacional]
    \footnotesize
    Todos los IC al 95\% contienen el valor de referencia.
    Error relativo $<1{,}2\%$ en las tres métricas.
    \textbf{El modelo está validado.}
  \end{resultadobox}
\end{frame}

\begin{frame}{Ejemplo 3 — Análisis de sensibilidad sobre $\lambda$}
  \begin{columns}[T]
    \begin{column}{0.46\textwidth}
      \begin{center}\scriptsize
        \renewcommand{\arraystretch}{1.25}
        \begin{tabular}{ccccc}
          \toprule
          $\lambda$ & $\rho$ & $W_q^{\text{P-K}}$ & $\bar W_q^{\text{sim}}$ & \textbf{Dir.} \\
          \midrule
          $6{,}4$  & $0{,}64$ & $0{,}53$ & $0{,}55$ & $\checkmark$ \\
          $7{,}2$  & $0{,}72$ & $0{,}77$ & $0{,}79$ & $\checkmark$ \\
          $8{,}0$  & $0{,}80$ & $1{,}20$ & $1{,}21$ & --- \\
          $8{,}8$  & $0{,}88$ & $2{,}20$ & $2{,}24$ & $\checkmark$ \\
          $9{,}6$  & $0{,}96$ & $7{,}20$ & $7{,}45$ & $\checkmark$ \\
          \bottomrule
        \end{tabular}
      \end{center}

      \vspace{4pt}
      \footnotesize
      \textbf{Tres lecturas:}
      \begin{itemize}\itemsep1pt
        \item Dirección correcta ($W_q$ aumenta con $\lambda$). $\checkmark$
        \item No-linealidad: $+20\%$ en $\lambda$ produce $+500\%$ en $W_q$
          — coherente con teoría de colas cerca de saturación. $\checkmark$
        \item Error relativo $<4\%$ en todos los puntos. $\checkmark$
      \end{itemize}
    \end{column}
    \begin{column}{0.50\textwidth}
      \begin{tikzpicture}
        \begin{axis}[width=\textwidth, height=6.0cm,
          title={$W_q$ vs.\ $\lambda$ (sensibilidad)},
          xlabel={$\lambda$ (pzas/h)}, ylabel={$W_q$ (h)},
          xmin=6, xmax=10, ymin=0, ymax=8,
          legend pos=north west]
          \addplot[thick, mobiblue, domain=6:9.8, samples=80]
            {x*0.06/(2*(1-x*0.2))};
          \addlegendentry{P-K analítica}
          \addplot[only marks, mark=*, mark size=2pt, princetonorange]
            coordinates {
              (6.4,0.55) (7.2,0.79) (8.0,1.21) (8.8,2.24) (9.6,7.45)
            };
          \addlegendentry{Simulación}
          \addplot[dashed, deepgreen, thick]
            coordinates {(8,0) (8,8)};
          \node[font=\tiny, deepgreen] at (axis cs:8.0,7.5) {base};
        \end{axis}
      \end{tikzpicture}
    \end{column}
  \end{columns}

  \vspace{2pt}
  \begin{resultadobox}[V\&V cerrada para el Ejemplo 3]
    \footnotesize
    Verificación: conservación, trazas y pruebas extremas correctas.
    Validación: $W_q$, $L_q$, $\rho$ caen dentro del IC 95\% de la referencia P-K
    (error rel.\ $<1{,}2\%$). Sensibilidad: dirección, magnitud y monotonía
    coherentes. \textbf{El modelo gana credibilidad para análisis fuera del
    alcance analítico} (percentiles, prioridades, múltiples colas).
  \end{resultadobox}
\end{frame}

% =============================================================
\section{Síntesis}
% =============================================================

\begin{frame}{El conjunto de evidencias}
  \footnotesize
  \textbf{Validar un modelo no es un único test, sino un \emph{conjunto} de
  evidencias:}
  \begin{center}
    \renewcommand{\arraystretch}{1.25}
    \begin{tabular}{p{4.2cm} p{4.2cm} p{4.4cm}}
      \toprule
      \textbf{Sesión 1: Verificación} & \textbf{Sesión 2: Validación}
        & \textbf{Aporta evidencia sobre} \\
      \midrule
      \emph{Debugging}, aserciones    & --- & Código sin defectos puntuales \\
      Trazas + invariantes            & --- & Código respeta lógica del modelo \\
      Pruebas unitarias               & --- & Casos límite y analíticos correctos \\
      \midrule
      --- & \emph{Face validity}      & Modelo razonable a juicio experto \\
      --- & Validación estructural    & Lógica del modelo refleja el sistema \\
      --- & $\chi^2$, KS sobre datos  & Distribuciones de \emph{inputs} correctas \\
      --- & Calibración (división de muestra)
                                      & Parámetros ajustados generalizan \\
      --- & I/O con $t$ pareada       & Salidas del modelo $\approx$ realidad \\
      --- & Test de Turing            & Salidas indistinguibles para expertos \\
      --- & Análisis de sensibilidad  & Modelo responde \emph{como debe} a cambios \\
      \bottomrule
    \end{tabular}
  \end{center}

  \vspace{4pt}
  \begin{notabox}[Definición operativa de modelo válido]
    \footnotesize
    Un modelo es ``válido'' cuando la evidencia acumulada en estos planos es
    suficiente para que las decisiones que se tomen con él tengan
    \textbf{respaldo razonable}.
  \end{notabox}
\end{frame}

\begin{frame}{Errores frecuentes en V\&V}
  \begin{center}\footnotesize
    \renewcommand{\arraystretch}{1.3}
    \begin{tabular}{p{4.4cm}p{4.4cm}p{4.6cm}}
      \toprule
      \textbf{Error} & \textbf{Síntoma} & \textbf{Cómo evitarlo} \\
      \midrule
      Validar con los mismos datos usados para calibrar
        & Error relativo $<1\%$ ``demasiado bueno''
        & Partir el histórico en $\sim 70/30$ — calibrar / validar \\
      Validar solo la media
        & La media coincide pero los percentiles no
        & Comparar también $W_q^{(95)}$, varianza y forma de la distribución \\
      Olvidar el período de calentamiento
        & Sesgo persistente en los IC
        & Descartar el transiente (ver T3.4) \\
      Usar pocas réplicas
        & IC muy ancho, no decidible
        & Aumentar $n_R$ hasta que el IC sea $<5\%$ del valor estimado \\
      Validar solo en el punto base
        & Modelo válido localmente, falso fuera
        & Análisis de sensibilidad con varios escenarios \\
      Confundir significancia estadística con práctica
        & ``Rechazar $H_0$'' por diferencia trivial
        & Acompañar la prueba con error relativo y juicio del experto \\
      \bottomrule
    \end{tabular}
  \end{center}

  \vspace{6pt}
  \begin{alertabox}[Anti-patrón típico]
    \footnotesize
    ``El modelo coincide con los datos del año pasado, así que está validado para
    proyectar el próximo''. La validez \emph{histórica} no garantiza la validez
    \emph{prospectiva} si el sistema cambió (políticas, demanda, tecnología).
  \end{alertabox}
\end{frame}

\begin{frame}{Mensajes finales}
  \begin{enumerate}\small\itemsep4pt
    \item \textbf{Validación} (¿modelo correcto?) es \textbf{conceptualmente
      primaria}; la verificación (¿correctamente construido?) es indispensable
      pero secundaria.
    \item Cada técnica formal vista hoy responde a \textbf{una pregunta distinta};
      no se sustituyen entre sí.
    \item Reportar siempre \textbf{intervalos de confianza}, no solo $p$-valores;
      distinguir \textbf{significancia estadística} de \textbf{relevancia operacional}.
    \item Calibrar y validar con \textbf{datos independientes} (división de muestra).
    \item El \textbf{análisis de sensibilidad} cierra el ciclo: dirección,
      magnitud y monotonía deben ser consistentes con la teoría.
    \item V\&V es \textbf{iterativa}: cada cambio en el modelo exige
      re-verificar y re-validar.
  \end{enumerate}

  \vspace{6pt}
  \begin{alertabox}[Mensaje para llevar]
    \footnotesize
    Un modelo verificado pero no validado es código limpio que responde la
    pregunta equivocada. Un modelo validado pero no verificado es una opinión
    correcta sobre un sistema mal programado. \textbf{Hacen falta los dos.}
  \end{alertabox}
\end{frame}

\begin{frame}{Próximo tema — T3.3: Análisis de salida (terminales)}
  \begin{columns}[T]
    \begin{column}{0.55\textwidth}
      \begin{situacionbox}[De la V\&V al uso del modelo]
        \footnotesize
        Con el modelo \emph{verificado} y \emph{validado}, podemos finalmente
        usarlo para responder preguntas que la teoría no resuelve.
        Pero antes de reportar \emph{cualquier} número, debemos saber
        \textbf{cómo promediar las réplicas}.
      \end{situacionbox}

      \vspace{4pt}
      \textbf{En T3.3 veremos:}
      \begin{itemize}\footnotesize\itemsep2pt
        \item Simulaciones \emph{terminales} vs.\ estacionarias.
        \item Réplicas independientes y construcción de IC.
        \item Determinación del número de réplicas necesarias.
        \item Reporte de métricas con incertidumbre.
      \end{itemize}
    \end{column}
    \begin{column}{0.42\textwidth}
      \begin{tikzpicture}[>=Stealth, node distance=1.4cm,
        proc/.style={rectangle, rounded corners=3pt, draw=mobiblue, thick,
          fill=mobilightblue, minimum width=2.4cm, minimum height=0.8cm,
          align=center, font=\scriptsize}]
        \node[proc, fill=mobilightgreen, draw=deepgreen] (in) {T3.1\\Entradas};
        \node[proc, below=of in, fill=mobilightgold, draw=mobigold] (vv)
          {\textbf{T3.2 — V\&V}\\(estamos aquí)};
        \node[proc, below=of vv] (ot) {T3.3\\Salida (terminal)};
        \node[proc, below=of ot] (oe) {T3.4\\Salida (estacionaria)};
        \draw[->, thick, mobiblue] (in) -- (vv);
        \draw[->, thick, mobiblue] (vv) -- (ot);
        \draw[->, thick, mobiblue] (ot) -- (oe);
      \end{tikzpicture}

      \vspace{4pt}
      \begin{notabox}[Continuidad]
        \footnotesize
        Lo que validamos aquí (medias, IC) es exactamente lo que T3.3 usará para
        \emph{reportar} los resultados de la simulación.
      \end{notabox}
    \end{column}
  \end{columns}
\end{frame}

% ------------------------------------------------------------
\begin{frame}{Referencias}
  \footnotesize
  \begin{thebibliography}{9}
  \bibitem{banks} Banks, J., Carson, J.\,S., Nelson, B.\,L., Nicol, D.\,M. (2010).
    \emph{Discrete-Event System Simulation} (5th ed.). Pearson. Caps.\ 1, 10.
  \bibitem{law} Law, A.\,M. (2015).
    \emph{Simulation Modeling and Analysis} (5th ed.). McGraw-Hill. Cap.\ 5.
  \bibitem{kelton} Kelton, W.\,D., Sadowski, R.\,P., Zupick, N.\,B. (2015).
    \emph{Simulation with Arena} (6th ed.). McGraw-Hill.
  \bibitem{ross} Ross, S.\,M. (2013).
    \emph{Simulation} (5th ed.). Academic Press.
  \bibitem{sargent} Sargent, R.\,G. (2013).
    Verification and Validation of Simulation Models.
    \emph{Journal of Simulation}, 7(1), 12--24.
  \bibitem{schruben} Schruben, L.\,W. (1980).
    Establishing the Credibility of Simulations.
    \emph{Simulation}, 34(3), 101--105.
  \bibitem{naylor} Naylor, T.\,H., Finger, J.\,M. (1967).
    Verification of Computer Simulation Models.
    \emph{Management Science}, 14(2), 92--101.
  \bibitem{balci} Balci, O. (1997).
    Verification, Validation and Accreditation of Simulation Models.
    \emph{Proceedings of the 1997 Winter Simulation Conference}.
  \bibitem{lilliefors} Lilliefors, H.\,W. (1969).
    On the KS Test for the Exponential Distribution with Mean Unknown.
    \emph{JASA}, 64(325), 387--389.
  \end{thebibliography}
\end{frame}

% ------------------------------------------------------------
\begin{frame}
  \begin{center}
    {\Huge \color{mobiblue}\textbf{¿Preguntas?}}\\[0.5cm]
    \large Verificación y Validación\\ en Simulación de Eventos Discretos\\[0.5cm]
    {\color{mobigold}\rule{0.5\textwidth}{1.2pt}}
  \end{center}
\end{frame}

\end{document}
